\subsection{\texorpdfstring{类型为 \(D_{l}\) 的代数 \((l\geq 2)\)}{类型为 D\_\{l\} 的代数 (l\textbackslash geq 2)}}

\begin{enumerate}
\def\labelenumi{(\Roman{enumi})}
\tightlist
\item
  设 V 为维数为 \(2l \geq 4\) 的向量空间，且 \(\varPsi\) 为 V 上指标为最大指标 l 的非退化对称双线性型。由《代数》第 IX 章，§4，第2节，V 可写成两个极大全迷向子空间 F 和 \(F'\) 的直和。令 \((e_i)_{1 \leq i \leq l}\) 为 F 的一个基，令 \((e_{-i})_{1 \leq i \leq l}\) 为 \(F'\) 的对偶基（关于由 \(F'\) 定义的 F 与 \(\varPsi\) 之间的对偶关系）。于是 \(e_1, \ldots, e_l, e_{-l}, \ldots, e_{-1}\) 是 V 的一个基；我们称其为 V 的 Witt 基。相对于此基，\(\varPsi\) 的矩阵是阶为 2l 的方阵 S，其元素全为 0，唯有位于第二对角线上的元素等于 1。代数 \(\mathfrak{g} = \mathfrak{o}(\varPsi)\) 可与方阵代数 \(\mathfrak{o}_S(2l, k)\) 识别，后者由满足 \(g = -S^t g S\) 的阶为 2l 的方阵 g 组成。其维数为 \(l(2l - 1)\)。直接计算表明，g 是如下形式的矩阵的集合
\end{enumerate}

\[
\left( \begin{array}{c c} A & B \\ C & D \end{array} \right)
\]

其中 A、B、C、D 是阶为 l 的方阵，并满足 \(B = -s^{t}Bs\)、\(C = -s^{t}Cs\) 和 \(D = -s^{t}As\)（s 是阶为 l 的矩阵，除位于第二对角线上的元素等于 1 外，其余元素全为 0）。

令 \(\mathfrak{h}\) 为属于 \(\mathfrak{g}\) 的对角矩阵的集合。这是 \(\mathfrak{g}\) 的交换子代数，其基由元素 \(H_{i} = E_{i,i} - E_{-i,-i}\)（\(1 \leq i \leq l\)）组成。令 \((\varepsilon_{i})\) 为与 \(\mathfrak{h}^{*}\) 对偶的 \((H_{i})\) 的基。对于 \(1 \leq i < j \leq l\)，置

\[
\left\{ \begin{array}{l l} X _ {\varepsilon_ {i} - \varepsilon_ {j}} & = E _ {i, j} - E _ {- j, - i} \\ X _ {- \varepsilon_ {i} + \varepsilon_ {j}} & = - E _ {j, i} + E _ {- i, - j} \\ X _ {\varepsilon_ {i} + \varepsilon_ {j}} & = E _ {i, - j} - E _ {j, - i} \\ X _ {- \varepsilon_ {i} - \varepsilon_ {j}} & = - E _ {- j, i} + E _ {- i, j}. \end{array} \right.\tag{11}
\]

这些元素构成 h 在 g 中的一个补空间的基。对于 \(h \in h\)，有

\[
[ h, X _ {\alpha} ] = \alpha (h) X _ {\alpha}
\]

对所有 \(\alpha \in \mathbb{R}\) 都成立，其中 \(\mathbb{R}\) 是由 \(\pm \varepsilon_i \pm \varepsilon_j\)（\(i < j\)）组成的集合。因此，\(\mathfrak{h}\) 是 \(\mathfrak{g}\) 的分裂嘉当子代数，且 \((\mathfrak{g}, \mathfrak{h})\) 的根是 \(\mathbb{R}\) 的元素。\(\mathbb{R}\) 的根系 \((\mathfrak{g}, \mathfrak{h})\) 在 \(\mathrm{D}_l\) 时为 \(l \geq 3\) 型，在 \(\mathrm{A}_1 \times \mathrm{A}_1\) 时为 \(\mathrm{D}_2\) 型（即 \(l = 2\) 型）（第 VI 章，§4，第8.I节，推广到 \(l = 2\) 的情形）。因此，\(\mathfrak{g}\) 是 \(\mathrm{D}_l\) 型的可分裂单李代数（当 \(l \geq 3\) 时）。

g 的每个分裂嘉当子代数都可由 g 的一个初等自同构变换为 h，因而也可由 \(\mathbf{O}(\varPsi)\) 的某个元素变换得到（参见~(VII)）；所以它是集合 \(h_{\beta}\)，其中的元素 g 相对于 V 的某个 Witt 基 \(\beta\) 的矩阵为对角矩阵。立即验证可知，在 \(h_{\beta}\) 下不变的唯一子空间是由 \(\beta\) 的某个子集生成的子空间。

由于代数 \(\mathfrak{o}_{S}(4,k)\) 和 \(\mathfrak{sl}(2,k)\times\mathfrak{sl}(2,k)\) 具有相同的根系，二者同构。同样，\(\mathfrak{o}_{S}(6,k)\) 和 \(\mathfrak{sl}(4,k)\) 同构（另参见~习题3）。下文设 \(l\geq3\)。

\begin{enumerate}
\def\labelenumi{(\Roman{enumi})}
\setcounter{enumi}{1}
\tightlist
\item
我们通过第 VI 章，§4，第8.V节确定 \(\mathbb{R}^{\vee}\)。我们得到
\end{enumerate}

\[
H _ {\varepsilon_ {i} - \varepsilon_ {j}} = H _ {i} - H _ {j}, \quad H _ {\varepsilon_ {i} + \varepsilon_ {j}} = H _ {i} + H _ {j}.
\]

\begin{enumerate}
\def\labelenumi{(\Roman{enumi})}
\setcounter{enumi}{2}
\tightlist
\item
置 \(\alpha_{1} = \varepsilon_{1} - \varepsilon_{2},\alpha_{2} = \varepsilon_{2} - \varepsilon_{3},\ldots ,\alpha_{l - 1} = \varepsilon_{l - 1} - \varepsilon_{l},\alpha_{l} = \varepsilon_{l - 1} + \varepsilon_{l}\)。由第 VI 章，§4，第8.II节，\((\alpha_{1},\dots,\alpha_{l})\) 是 R 的一个基 B；相对于 B 的正根是 \(\varepsilon_i\pm \varepsilon_j\)（\(i < j\)）。相应的 Borel 子代数 \(\mathfrak{b}\) 是属于 \(\mathfrak{g}\) 的上三角矩阵的集合。
\end{enumerate}

立即验证可知，在 b 下不变的非平凡向量子空间只有全迷向子空间 \(V_{1},\ldots,V_{l},V_{l}^{\prime}\)，其中 \(V_{i}\) 由 \(e_{1},\ldots,e_{i}\) 生成，\(V_{l}^{\prime}\) 由 \(e_{1},\ldots,e_{l-1},e_{-l}\) 生成，以及 \(V_{-1},\ldots,V_{-l+1}\) 的正交补 \(V_{1},\ldots,V_{l-1}\)；\(V_{-i}\) 的正交补 \(V_{i}\) 由 \(e_{1},\ldots,e_{l},e_{-l},\ldots,e_{-(i+1)}\) 生成。但是直接计算表明，若 \(a\in g\) 使 \(V_{l-1}\) 稳定，则其矩阵具有形式

\[
\left( \begin{array}{c c c} A & x & B \\ 0 & \left( \begin{array}{c c} \lambda & 0 \\ 0 & - \lambda \end{array} \right) & y \\ 0 & 0 & D \end{array} \right)
\]

其中 \(A, B, D\) 是阶为 \(l - 1\) 的方阵，\(x\)（分别 \(y\)）是具有 2 列和 \(l - 1\) 行（分别具有 2 行和 \(l - 1\) 列）的矩阵，且 \(\lambda \in k\)。由此 \(a\) 使 \(\mathrm{V}_l\) 和 \(\mathrm{V}_l'\) 稳定。因此，\(\mathfrak{b}\) 是 \(a \in \mathfrak{g}\) 中使迷向旗 \((\mathrm{V}_1, \ldots, \mathrm{V}_{l-1})\) 的所有元素稳定的集合。注意，前述事实和 Witt 定理（...）表明，唯一的相关极大全迷向子空间是 \(\mathrm{V}_l\) 和 \(\mathrm{V}_l'\)，均包含 \(\mathrm{V}_{l-1}\)。

我们称一个迷向旗为拟极大迷向旗，如果它由 l-1 个维数为 \(1,\ldots,l-1\) 的全迷向子空间组成。于是，如第2节中那样可知，对于任意拟极大迷向旗 \(\delta\)，使 \(b_{\delta}\) 的各元素稳定的 \(a\in g\) 的集合 \(\delta\) 是 g 的一个 Borel 子代数，并且映射 \(\delta\mapsto b_{\delta}\) 是从拟极大迷向旗集合到 Borel 子代数集合的双射。

我们称一个迷向旗为真的，如果它不同时含有一个 \(l\) 维子空间和一个 \(l-1\) 维子空间。令 \(\delta\) 为这样的迷向旗，并令 \(\mathfrak{p}_{\delta}\) 为由 \(a \in \mathfrak{g}\) 中使 \(\delta\) 的各元素稳定的元素组成的集合。若它包含于 \(\delta \subset \{\mathrm{V}_1, \ldots, \mathrm{V}_l, \mathrm{V}_l'\}\)，则 \(\mathfrak{p}_{\delta}\) 是 \(\mathfrak{g}\) 的一个抛物子代数，包含 \(\mathfrak{b}\)，并且容易验证，唯一不等于 \(\neq \{0\}\) 且在 \(\mathfrak{p}_{\delta}\) 下稳定的全迷向子空间是 \(\delta\) 的元素。由于每类有 \(2^{l-2}\) 个 \(\{\mathrm{V}_1, \ldots, \mathrm{V}_l, \mathrm{V}_l'\}\) 中的真的迷向旗，其中包含 \(\mathrm{V}_{l-1}\)（分别包含 \(\mathrm{V}_l\)、包含 \(\mathrm{V}_l'\)、既不包含 \(\mathrm{V}_{l-1}\)、也不包含 \(\mathrm{V}_l\) 和 \(\mathrm{V}_l'\)），这给出了 \(2^l\) 个抛物子代数。由此，如上所述，映射 \(\mathfrak{b}\) 是从真的迷向旗集合到 \(\delta \mapsto \mathfrak{p}_{\delta}\) 的抛物子代数集合（属于 \(\mathfrak{g}\)）的双射。

\begin{enumerate}
\def\labelenumi{(\Roman{enumi})}
\setcounter{enumi}{3}
\tightlist
\item
与 \(\alpha_{1},\ldots,\alpha_{l}\) 对应的基本权由第 VI 章，§4，第8.VI节给出：
\end{enumerate}

\[
\varpi_ {i} = \varepsilon_ {1} + \varepsilon_ {2} + \dots + \varepsilon_ {i} (1 \leq i \leq l - 2)
\]

\[
\begin{array}{c} \varpi_ {l - 1} = \frac {1}{2} (\varepsilon_ {1} + \varepsilon_ {2} + \ldots + \varepsilon_ {l - 2} + \varepsilon_ {l - 1} - \varepsilon_ {l}) \\ \varpi_ {l} = \frac {1}{2} (\varepsilon_ {1} + \varepsilon_ {2} + \ldots + \varepsilon_ {l - 2} + \varepsilon_ {l - 1} + \varepsilon_ {l}). \end{array}
\]

令 \(\sigma\) 为 \(\mathfrak{g}\) 在 V 上的恒等表示。外幂 \(\bigwedge^{r}\sigma\) 作用于 \(\mathrm{E} = \bigwedge^{r}(\mathrm{V})\)。若 \(h\in \mathfrak{h}\)，有

\[
\sigma (h) e _ {i} = \varepsilon_ {i} (h) e _ {i}, \quad \sigma (h) e _ {- i} = - \varepsilon_ {i} (h) e _ {- i}
\]

对于 \(1 \leq i \leq l\)，由此，对于 \(1 \leq r \leq l\)，\(\varepsilon_{1} + \cdots + \varepsilon_{r}\) 是 \(\bigwedge^{r} \sigma\) 的最高权，权为 \(\varepsilon_{1} + \cdots + \varepsilon_{r}\) 的元素正是与 \(e_{1} \wedge \cdots \wedge e_{r}\) 成比例的元素。

我们将证明，当 \(1 \leq r \leq l - 2\) 时，表示 \(\bigwedge^{r} \sigma\) 是权为 \(\varpi_{r}\) 的基本表示。

为此，只需证明 \(\bigwedge^{r}\sigma\) 在 \(1 \leq r \leq l - 1\) 时并非不可约（注意表示 \(\bigwedge^{l}\sigma\) 不可约，参见~习题10），或者证明 \(T_{r}\) 中包含 \(\bigwedge^{r}V\) 且在 g 下稳定的最小子空间 \(e_{1} \wedge \cdots \wedge e_{r}\) 就是整个 \(\bigwedge^{r}V\)。当 r = 1 时这立即成立。当 r = 2 时，如第2节中那样，我们看到 \(\bigwedge^{2}\sigma\) 等价于 g 的伴随表示；由于 g 是单的，该表示不可约。通过像第2节中那样对 l 作归纳，并设 \(l - 1 \geq r \geq 3\)，即可完成证明。

现在确定最高权为 \(\varpi_{l-1}\) 和 \(\varpi_{l}\) 的基本表示。令 Q 为二次型 \(x \mapsto \frac{1}{2}\Psi(x, x)\)。我们在第2.IV节中定义了 Clifford 代数 \(\lambda\) 在 \(\mathrm{C}(\mathrm{Q})\) 上的旋量表示 \(N = \bigwedge F'\)。立即验证可知，N 的子空间 \(N_{+}\)（分别为 \(N_{-}\)）由 p 为偶数（分别为奇数）时的 \(\bigwedge^{p}F'\) 之和给出，并且在 \(\lambda\) 对 \(\mathrm{C}^{+}(\mathrm{Q})\) 的限制下稳定。因此，\(\lambda_{+}\) 和 \(\lambda_{-}\) 分别是 \(\mathrm{C}^{+}(\mathrm{Q})\) 在 \(N_{+}\) 和 \(N_{-}\) 上的表示，也是 \(\mathrm{C}^{+}(\mathrm{Q})\) 的半旋量表示（《代数》第 IX 章，§9，第4节）；它们不可约，维数为 \(2^{l-1}\)，且不等价。令 \(\rho_{+} = \lambda_{+} \circ f\) 和 \(\rho_{-} = \lambda_{-} \circ f\) 为相应的 g 的不可约表示（第2节，引理1 (vi)）。根据引理1 (i)，有

\[
f (H _ {i}) = \frac {1}{2} (e _ {i} e _ {- i} - e _ {- i} e _ {i}) = e _ {i} e _ {- i} - \frac {1}{2} = \frac {1}{2} - e _ {- i} e _ {i}
\]

并且如第2.IV节中那样可见，对于 \(h \in h\) 和 \(1 \leq i_{1} < \cdots < i_{k} \leq l\)，

\[
\begin{array}{l} \lambda \circ f (h) (e _ {- i _ {1}} \wedge \dots \wedge e _ {- i _ {k}}) \\ = (\frac {1}{2} (\varepsilon_ {1} + \dots + \varepsilon_ {l}) - (\varepsilon_ {i _ {1}} + \dots + \varepsilon_ {i _ {k}})) (h) (e _ {- i _ {1}} \wedge \dots \wedge e _ {- i _ {k}}). \end{array}
\]

因此，\(\rho_{+}\)（分别 \(\rho_{-}\)）的最高权是 \(\varpi_{l}\)（分别 \(\varpi_{l-1}\)）。

我们称 \(\rho_{+}\) 和 \(\rho_{-}\) 为 g 的半旋量表示。它们的所有权都是单权。我们还称 \(\rho = \lambda \circ f = \rho_{+} \oplus \rho_{-}\) 为 g 的旋量表示。

\begin{enumerate}
\def\labelenumi{(\Alph{enumi})}
\setcounter{enumi}{21}
\tightlist
\item
对于 \(1 \leq r \leq l - 2\)，基本表示 \(\bigwedge^{r} \sigma\) 是正交型的：它保持 \(\varPsi\) 在 \(\bigwedge^{r} \mathrm{V}\) 上的延拓不变。
\end{enumerate}

现在考虑 \(\rho\) 的旋量表示 \(\mathfrak{g}\)。如第2节中那样，我们说明，对于 \(1 \leq i, j \leq l, i \neq j\)，

\[
\begin{array}{r} f (X _ {\varepsilon_ {i} - \varepsilon_ {j}}) = \pm e _ {i} e _ {- j} \\ f (X _ {\varepsilon_ {i} + \varepsilon_ {j}}) = \pm e _ {i} e _ {j} \\ f (X _ {- \varepsilon_ {i} - \varepsilon_ {j}}) = \pm e _ {- i} e _ {- j}. \end{array}
\]

由此，在第2.V节中引入的非退化双线性型 \(\varPhi\) 在 \(\rho(\mathfrak{g})\) 下不变。因此，旋量表示 \(\rho\) 保持一个非退化型不变；当 \(l \equiv 0, -1 \pmod{4}\) 时该型对称，当 \(l \equiv 1, 2 \pmod{4}\) 时该型交错。

当 \(l\) 为奇数时，\(\varPhi\) 交换 \(\mathrm{N}_{+}\) 与 \(\mathrm{N}_{-}\)；当 \(l \equiv 0 \pmod{4}\) 时，它具有相应的性质；当 \(l \equiv 2 \pmod{4}\) 时，\(w_{0} = -1\) 具有所述作用。

另一方面，若 \(l\) 为奇数，则 \(\mathrm{N}_{+}\) 和 \(\mathrm{N}_{-}\) 关于 \(\varPhi\) 都具有相应的性质。此外，\(-w_{0}(\alpha_{i}) = \alpha_{i}\) 对 \(1 \leq i \leq l - 2\) 成立，\(-w_{0}(\alpha_{l}) = \alpha_{l - 1}\)，而 \(-w_{0}(\alpha_{l - 1}) = \alpha_{l}\)，所以 \(-w_{0}(\varpi_{l}) = \varpi_{l - 1}\)。

\begin{enumerate}
\def\labelenumi{(\Roman{enumi})}
\setcounter{enumi}{5}
\tightlist
\item
对于所有 \(x \in \mathfrak{g}\)，\(\sigma(x)\) 的特征多项式具有形式
\end{enumerate}

\[
\mathrm{T} ^ {2 l} + f _ {1} (x) \mathrm{T} ^ {2 l - 1} + \dots + f _ {2 l} (x).
\]

我们像第3.VI节中那样看到 \(f_{1}=f_{3}=f_{5}=\cdots=0\)。由第 VI 章，§4，第8.IX节和§8，第3节，定理1，存在 g 上的多项式函数 \(\tilde{f}\)，使得 \(f_{2}, f_{4}, \ldots, f_{2l-2}, \tilde{f}\) 生成 \(\mathrm{I}(\mathfrak{g}^{*})\)（即 g 上不变多项式函数的代数），它们代数无关，并且还有 \(\tilde{f}^{2}=(-1)^{l}f_{2l}\)。

对于 \(x \in \mathfrak{g}\)，有 \({}^t(Sx) = {}^t xS = -Sx\)，所以可以考虑 \(\operatorname{Pf}(Sx)\)，它是 \(x\) 的多项式函数。现在

\[
f _ {2 l} (x) = \det (x) = (- 1) ^ {l} \det (S x) = (- 1) ^ {l} (\operatorname{Pf} (S x)) ^ {2}.
\]

因此，可取 \(\tilde{f}(x) = \mathrm{Pf}(Sx)\)。

\begin{enumerate}
\def\labelenumi{(\Roman{enumi})}
\setcounter{enumi}{6}
\tightlist
\item
回顾（§5，第3节，命题5的推论1），\(\mathrm{Aut}(\mathfrak{g}) / \mathrm{Aut}_0(\mathfrak{g})\) 可与 \(\mathrm{Aut}(\mathrm{D})\) 识别，后者是 \((\mathfrak{g},\mathfrak{h})\) 的邓金图 D 的自同构群。当 \(l\neq 4\) 时，\(\mathrm{Aut}(\mathrm{D})\) 是阶为 2 的群，由\(\alpha_{1},\ldots ,\alpha_{l}\)的置换组成，这些置换保持\(\alpha_{1},\ldots ,\alpha_{l - 2}\)不变。当 \(l = 4\) 时，\(\mathrm{Aut}(\mathrm{D})\) 由保持 \(\alpha_{1},\ldots ,\alpha_{4}\) 不变的 \(\alpha_{2}\) 的置换组成；它同构于 \(\mathfrak{S}_3\)（参见~第 VI 章，§4，第8.XI节）。在所有情形下，\(\mathrm{Aut}(\mathrm{D})\) 中保持 \(\alpha_{1}\) 不变的元素组成的子群的阶为 2。我们将相应的 \(\operatorname{Aut}'(\mathfrak{g})\) 的子群记为 \(\operatorname{Aut}(\mathfrak{g})\)；当 \(\operatorname{Aut}'(\mathfrak{g}) = \operatorname{Aut}(\mathfrak{g})\) 时，\(l \neq 4\)，当 l = 4 时 \((\operatorname{Aut}(\mathfrak{g}) : \operatorname{Aut}'(\mathfrak{g})) = 3\)；此外，
\end{enumerate}

\[
\left(\operatorname{Aut} ^ {\prime} (\mathfrak {g}): \operatorname{Aut} _ {0} (\mathfrak {g})\right) = 2.
\]

元素 \(s \in \operatorname{Aut}(\mathfrak{g})\) 属于 \(\operatorname{Aut}'(\mathfrak{g})\)，当且仅当 \(\sigma \circ s\) 与 \(\sigma\) 等价（这是因为 \(\varpi_{1}\) 是 \(\sigma\) 的最高权）。如第2.VII节中那样，我们得到 \(\operatorname{Aut}'(\mathfrak{g})\) 可与 \(\Sigma/k^{*}\) 识别，其中 \(\Sigma\) 是关于 \(\Psi\) 的 V 的相似变换群。

令 \(s \in \Sigma\)，令 \(\lambda(s)\) 为 s 的乘子。如果 \(s\) 是正相似，则 \(\det(s) = \lambda(s)^l\)；如果 \(s\) 是逆相似，则 \(\det(s) = -\lambda(s)^l\)。因此，\(s\) 也是正相似；正相似构成 \(\Sigma_0\) 的一个指数为 2 的子群，并且 \(\Sigma\) 满足 \(\Sigma_0 \supset k^*\)。群 \(\Sigma_0 / k^*\) 等于 \(\operatorname{Aut}_0(\mathfrak{g})\)（它是 \(\operatorname{Aut}'(\mathfrak{g}) = \Sigma / k^*\) 的子群）。事实上，只需在 \(k\) 代数闭时验证这一点：此时 \(\operatorname{Aut}_0(\mathfrak{g}) = \operatorname{Aut}_e(\mathfrak{g})\) 等于其换位子群（...），所以包含于 \(\Sigma_0 / k^*\)；而二者在 \(\Sigma / k^*\) 中的指数都为 2，故二者相等。

另一方面，正如第3.VII节中那样，有一个正合序列

\[
1 \longrightarrow \mathbf {S O} (\varPsi) / \{1, - 1 \} \longrightarrow \varSigma_ {0} / k ^ {*} \longrightarrow k ^ {*} / k ^ {* 2} \longrightarrow 1.
\]

将 \(\mathbf{SO}(\varPsi)/\{1,-1\}\) 识别为 \(\Sigma_{0}/k^{*}=\mathrm{Aut}_{0}(\mathfrak{g})\) 的一个子群。由于 \(k^{*}/k^{*2}\) 是交换群，有 \(\mathrm{Aut}_{e}(\mathfrak{g})\subset\mathbf{SO}(\varPsi)/\{1,-1\}\)。事实上，可以证明（习题11），\(\mathrm{Aut}_{e}(\mathfrak{g})\) 等于 \(\mathbf{SO}(\varPsi)/\{1,-1\}\) 的既约正交群 \(\mathbf{O}_{0}^{+}(\varPsi)\) 在 \(\varPsi\) 中的像（《代数》第 IX 章，§9，第5节）。

\begin{enumerate}
\def\labelenumi{(\Roman{enumi})}
\setcounter{enumi}{7}
\tightlist
\item
典则双线性型 \(\varPhi_{\mathrm{R}}\) 在 \(\mathfrak{h}^*\) 上为
\end{enumerate}

\[
\Phi_ {\mathrm{R}} \left(\xi_ {1} \varepsilon_ {1} + \dots + \xi_ {l} \varepsilon_ {l}, \xi_ {1} ^ {\prime} \varepsilon_ {1} + \dots + \xi_ {l} ^ {\prime} \varepsilon_ {l}\right) = \frac {1}{4 (l - 1)} \left(\xi_ {1} \xi_ {1} ^ {\prime} + \dots + \xi_ {l} \xi_ {l} ^ {\prime}\right)
\]

（第 VI 章，§4，第8.V节）。因此，Killing 型在 \(\mathfrak{h}\) 上的限制为

\[
\Phi \left(\xi_ {1} H _ {1} + \dots + \xi_ {l} H _ {l}, \xi_ {1} ^ {\prime} H _ {1} + \dots + \xi_ {l} ^ {\prime} H _ {l}\right) = 4 (l - 1) \left(\xi_ {1} \xi_ {1} ^ {\prime} + \dots + \xi_ {l} \xi_ {l} ^ {\prime}\right).
\]

\begin{enumerate}
\def\labelenumi{(\Roman{enumi})}
\setcounter{enumi}{8}
\tightlist
\item
回顾由公式（11）定义的 \(X_{\alpha}\)（\(\alpha \in \mathbb{R}\)）。容易验证，对于 \([X_{\alpha}, X_{-\alpha}] = -H_{\alpha}\)，有 \(\alpha \in \mathbb{R}\)。另一方面，映射 \(\theta : a \mapsto -^{t}a\) 是 \(\mathfrak{g}\) 的自同构，且对所有 \(\theta(X_{\alpha}) = X_{-\alpha}\) 有 \(\alpha \in \mathbb{R}\)。因此，\((X_{\alpha})_{\alpha \in \mathbb{R}}\) 是 \((\mathfrak{g}, \mathfrak{h})\) 中的谢瓦莱系。
\end{enumerate}

假设 \(k = \mathbf{Q}\)。嘉当子代数 \(\mathfrak{h}\) 在 \(l\) 为奇数时有三个许可格，在 \(l\) 为偶数时有四个许可格。特别地，格 \(\mathcal{H}\) 由 \(H_i\) 生成，且是 \(\mathfrak{g}\) 的许可格。因此，\(\mathfrak{o}_S(2l,\mathbf{Z})\) 是谢瓦莱序 \(\mathcal{H} + \sum \mathbf{Z}.X_\alpha\)，属于 \(\mathfrak{g}\)。由于 \(X_{\alpha}^{2} = 0\) 对所有 \(\alpha \in \mathrm{R}\) 都有，可知格 \(\mathcal{V}\) 位于 \(\mathrm{V}\) 中，由 Witt 基 \((e_i)\) 生成；它是 \(\mathrm{V}\) 中关于 \(\mathfrak{o}_S(2l,\mathbf{Z})\) 的容许格。\(\bigwedge^r\mathcal{V}\) 在 \(\bigwedge^r\mathrm{V}\) 中也如此。

另一方面，若取 \(\mathrm{P}(\mathrm{R}^{\vee}) = \mathbf{Z}\) 作为许可格，并取 \(\frac{1}{2} \sum_{i=1}^{l} H_i + \mathcal{H}\) 作为谢瓦莱序，则可知 \(\mathcal{G} = \mathrm{P}(\mathrm{R}^{\vee}) + \sum \mathbf{Z}. X_{\alpha}\) 中的 \(\bigwedge^r \mathrm{V}\) 只有在 \(r\) 为偶数时才有容许格；此时 \(\bigwedge^r \mathcal{V}\) 是容许格。

考虑约化李代数 \(\mathfrak{o}(\varPsi) + \mathbf{Q}.1\)；立即可见，格 \(\tilde{\mathcal{G}} = (\mathfrak{o}(\varPsi) + \mathbf{Q}.1) \cap \mathfrak{gl}(2l, \mathbf{Z})\) 是一个谢瓦莱序。谢瓦莱序 \(\mathcal{G}\) 是 \(\tilde{\mathcal{G}}\) 平行于中心 \(\mathfrak{o}(\varPsi)\) 投影到 \(\mathbf{Q}.1\) 上所得的序。

最后，如第2节中那样可见，由 \(N_{+}\)（分别由 \(N_{-}\)）生成的格 \(e_{-i_{1}}\wedge\cdots\wedge e_{-i_{2k}}\)（分别 \(e_{-i_{1}}\wedge\cdots\wedge e_{-i_{2k+1}}\)）是谢瓦莱序 \(\mathrm{Q}(\mathrm{R}^{\vee})+\sum_{\alpha\in\mathbb{R}}\mathbf{Z}.X_{\alpha}\) 的半旋量表示的容许格。另一方面，\(N_{+}\) 和 \(N_{-}\) 对 \(\mathfrak{o}_{S}(2l,\mathbf{Z})\) 没有容许格。

